\begin{table}[htbp]
\centering
\begin{threeparttable}
\caption{Matched comparisons of information, functional form, and forecast combination.}
\label{tab:matched-comparisons}
\scriptsize
\setlength{\tabcolsep}{4pt}
\begin{tabular}{@{}llrrr@{}}
\toprule
Comparison & Horizon & S\&P~500 & FTSE~100 & DAX \\
\midrule
\multicolumn{5}{l}{\textit{Panel A: Mean QLIKE differences}} \\
@@ROW_0@@
@@ROW_1@@
@@ROW_2@@
@@ROW_3@@
@@ROW_4@@
@@ROW_5@@
@@ROW_6@@
\addlinespace
\multicolumn{5}{l}{\textit{Panel B: Mean negative log score differences}} \\
@@ROW_7@@
@@ROW_8@@
\bottomrule
\end{tabular}
\begin{tablenotes}[flushleft]
\footnotesize
\item Notes: Entries are mean loss differences, with two-sided DM $p$-values in brackets; the tests are descriptive and unadjusted for multiplicity. Realised information is RV-NN-SV minus NN-SV; neural form is RV-NN-SV minus RV-LIN-SV; each combination is the corresponding ensemble minus RV-NN-SV. Negative differences favour the first model. QLIKE common-sample sizes for $\ell=(1,5,10)$ are $(1012,965,910)$ for the S\&P~500, $(1036,1024,1009)$ for the FTSE~100, and $(1029,1005,975)$ for the DAX. Neural-form sample sizes are 921, 994, and 1012. The one-day negative log score samples match the one-day QLIKE samples.
\end{tablenotes}
\end{threeparttable}
\end{table}
